% add the answers flag to print the solutions.
\documentclass[addpoints]{exam}

\usepackage[utf8]{inputenc}
\usepackage{calc}
\usepackage[dvipsnames]{xcolor}
\usepackage{ifthen}
\usepackage{tcolorbox}

% \makebox[\textwidth]{Name:\enspace\hrulefill}
%  
% \vspace{5mm}
%\makebox[\textwidth]{Instructor’s name:\enspace\hrulefill}

\definecolor{MyLightGray}{gray}{0.60}

\newcommand{\truefalse}
{		\begin{checkboxes}
			 \choice True
			 \choice False. Explanation: 
		\end{checkboxes}
		\vspace{10pt}
}


% Check if 'answers' is an option
\newcommand{\ifanswers}[2]{%
  \ifthenelse{\boolean{exam@answers}}{#1}{#2}%
}



\newcommand{\answerspace}[1]
{
	\ifprintanswers
	\else
		\setlength{\mylength}{10pt}%
		% \vspace{#1\mylength}
		% Make the grey box
		\par\nobreak\noindent
		\makebox[\linewidth][l]{%
			% set inner padding to zero, so only the frame shows
			{\setlength{\fboxsep}{0pt}%
			% draw a color frame, white (or transparent) inside
			\fcolorbox{MyLightGray}{white}{%
			% the inner box has zero width content but fixed height
			\parbox[t][#1\mylength][c]{%
				% subtract twice the frame‐thickness from the full width
				\dimexpr\linewidth-2\fboxrule\relax
			}{\strut}%
			}%
			}%
		}%
		\par
		\vspace{3\mylength}

  	\fi
}
\newlength{\mylength}

\newcommand{\code}[1]{{\texttt{#1}}}

\newcommand{\answer}[1]{\begin{solution}{#1}\end{solution}}




% For those sweet, sweet acmart fonts.
\usepackage[T1]{fontenc}
\usepackage[tt=false, type1=true]{libertine}
\usepackage[varqu]{zi4}                  % Inconsolata for \ttfamily

%\usepackage{listings}
\usepackage{color}
\usepackage{lastpage}
\usepackage{verbatim}
\usepackage{fancyvrb}
\usepackage{graphicx}
\usepackage{float}
\usepackage{minted}
\usepackage{multicol}

\newcommand{\includecode}[2][c]{\lstinputlisting[caption=, escapechar=, language=#1]{#2}}
\RecustomVerbatimCommand{\VerbatimInput}{VerbatimInput}{}

\newcommand{\Code}[1]{\mintinline{python}{#1}}

\renewcommand{\solutiontitle}{\noindent\textbf{Solution sketch: }}

\definecolor{SolutionColor}{rgb}{0.8,0.9,1}

\cfoot{Page \thepage\ of \pageref{LastPage}}

\rhead{Initials: \rule{20mm}{0.4pt}}


%%%%%%%%%%%%%%%%%%%%%%%%%%%%%%%%%%%%%%%%
\begin{document}
\fvset{numbers=left,numbersep=3pt}
\thispagestyle{empty}

\begin{center}
	{\LARGE CS 408 Fall 2026 -- Practice Midterm 2 \\}
\vspace{.5in}


\fbox{\fbox{\parbox{5.5in}{\centering
	\begin{itemize}
        \item We have deliberately made this practice midterm harder,
          more open-ended, and longer than the actual midterm, since
          there are no time constraints (and the course staff is not
          on the hook for grading your answers :)). \textbf{The goal
            here is to help you study}: if you can ace the practice
          midterm, you should have no problem with the real
          deal. Thus, while we will provide you with solutions next
          week, you should try to answer all these questions on your
          own, before consulting our answers.

        \item Try to be concise: when a question asks for an
          explanation, a few sentences should suffice. Aim for
          quality, not quantity!

        \item Show enough of your reasoning so that we can award
          partial credit for a correct approach.

        \item Answer the questions in the spaces provided. Your answer
          does not need to take up all the space provided (hopefully
          we've given you more than enough room!). If you do run out
          of space, make a note and then use the empty pages at the
          back of the exam to complete your answer. Be sure to clearly
          label any answers on those pages with the corresponding
          question.

        \item Several questions ask you to provide something ``if
          possible.'' If it is not possible, you should plainly say so
          and then explain why that is the case.

        \item Coding questions should be answered in Python, unless
          stated otherwise.

        \item Unless otherwise noted, CFG nodes represent basic
          blocks. When necessary, refer to nodes using the labels
          given in the figure.

        \item Please double check that your copy of the exam
          is complete; you should have \pageref{LastPage}
          pages.

        \item You have 75 minutes to complete the exam.
        \end{itemize}
	}}}
\end{center}

\vspace{5mm}
\begin{center}
	\makebox[0.6\textwidth]{
		\combinedgradetable[h][questions]
	}
\end{center}

\begin{questions}

\newpage

\question Short answers.
\begin{parts}

% --- from the oracles / PBT material ---------------------
\part[2] You test the property \Code{reverse(reverse(xs)) == xs} on a
  thousand generated lists and it passes every time. Can you conclude
  that \Code{reverse} is correct? Explain briefly. %
  \begin{solutionorbox}[1.5in]
    No. An implementation that returns its argument unchanged
    satisfies the round-trip property perfectly, as does one that
    rotates the list. The property only captures part of
    \Code{reverse}'s functionality.
  \end{solutionorbox}

% --- synthesis across the whole term ---------------------
\part[2] Order the following by how much information about the program
  under test they require, from least to most: (A) symbolic execution,
  (B) black-box fuzzing, (C) coverage-guided grey-box fuzzing. %
  \begin{solutionorbox}[1.5in]
    B $<$ C $<$ A. Black-box fuzzing needs only the ability to run the
    program; grey-box adds lightweight coverage instrumentation;
    symbolic execution needs the source (or bytecode) in order to build
    path conditions at all.
  \end{solutionorbox}
\end{parts}

% ========================================================
% Test Oracles
% ========================================================

\question[3] Identify three categories of test oracles that we have
  discussed in class, and list the relative strengths and weaknesses of
  each.
  \answerspace{24} %
  \begin{solution} %
    \begin{description}
    \item[Implicit oracles] \emph{Pros:} broadly applicable, require no
      extra input from the user, easy to automate. \emph{Cons:} weak
      correctness guarantees --- they only catch crashes and the like.
    \item[Human oracles] \emph{Pros:} very flexible; applicable when
      correctness is hard to specify (think UI). \emph{Cons:} expensive
      and slow; does not scale.
    \item[Reference oracles] \emph{Pros:} strong correctness
      guarantees, fully automatable, ideal for regression testing.
      \emph{Cons:} requires a trusted reference implementation, which
      may not exist.
    \item[Property-based oracles] \emph{Pros:} apply to every input,
      no reference needed. \emph{Cons:} properties may be too weak to
      pin down correctness.
    \end{description}
  \end{solution}

  \newpage
% ========================================================
% Meta-morphic relations
% ========================================================

\question Consider an image classifier that reports \textbf{the number
    of cats in an image}.

\begin{parts}

\part[3] Give two metamorphic relations for this classifier that could
  reveal a fault. At least one should explore a way the output of the
  classifier should \emph{change}, rather than requiring it to stay
  the same. %
  \answerspace{16} %
  \begin{solution} %
    Many correct answers. Two of different kinds:

    \textbf{Invariance:} rotate the image, or brighten it slightly
    $\longrightarrow$ the reported count is unchanged.

    \textbf{Change:} paste a second copy of the image beside the first
    $\longrightarrow$ the count doubles. Or: paste in a picture of one
    more cat $\longrightarrow$ the count goes up by exactly one.
  \end{solution}

\part[3] Give an example of a faulty classifier that satisfies one of
  the relations you gave above but is still obviously wrong, and
  explain why the fault evaded detection. %
  \answerspace{14} %
  \begin{solution} %
    A classifier that always returns a constant, e.g., \Code{0}
    satisfies every invariance relation: rotating the image does not
    change the answer, because nothing changes the answer.

    This is the general trap with invariance relations. They constrain
    only how outputs relate, so a function that ignores its input
    entirely is maximally consistent. The relations that pin down how
    the output must \emph{change} are much harder to satisfy by
    accident, which is why the previous part asked for one of each.
  \end{solution}

\part[2] In one or two sentences, explain what this tells us about the
  pitfualls of using metamorphic relations as the \emph{only} test
  oracle? %
  \answerspace{14} %
  \begin{solution} %
    A metamorphic relation never establishes that an answer is
    correct, only that the behavior of the SUT is internally
    consistent, and even a faulty system can give answers that are
    perfectly consistent.
  \end{solution}

\end{parts}

% ========================================================
% Test Input Generators
% ========================================================

Here is a Hypothesis strategy that produces binary trees, where a tree
is represented as either \Code{None}, for an empty tree, or a triple
\Code{(left, value, right)}.:

\begin{minted}{python}
@st.composite
def trees(draw, lo=-1000, hi=1000, max_depth=5):
    if max_depth <= 0 or draw(st.integers(0, 2)) == 0:
        return None
    return (draw(trees(lo, hi, max_depth - 1)),
            draw(st.integers(min_value=lo, max_value=hi)),
            draw(trees(lo, hi, max_depth - 1)))
\end{minted}

\begin{parts}

\part[3] Suppose that we want to use this strategy to test a function,
  \Code{is_member(t, v)}, that searches a binary search tree for a
  particular element. One way to do so would be to use \Code{trees} as
  is, filtering out any randomly generated trees that are not valid
  BSTs, e.g., using \Code{.filter(is_bst)}, and then generating a
  random element to look for. In 3-4 sentences, explain why this is or
  is not an effective input generation strategy for testing
  \Code{is_member}.%
  \answerspace{14} %
  \begin{solution} %
    Almost none of the outputs produced by \Code{trees} will pass this
    filter: sampling 5000 trees using this strategy will produce about
    2800 with three or more nodes, and fewer than \textbf{1\%} of
    those are valid binary search trees.

    Moreover, among the trees that do get through, roughly
    \textbf{95\%} have one node or none, i.e., they are ordered by
    accident. So, the tests would pass while hardly ever exercising a
    tree of interesting size, and nothing in the output would tell us
    that.
  \end{solution}

\part[4] Modify \Code{trees} so that it produces \emph{only} binary
  search trees, with nothing rejected. You need not get the Hypothesis
  syntax exactly right, but be precise about what values each
  recursive call is allowed to produce. %
  \answerspace{18} %
  \begin{solution} %
    Draw the node's value first, then \textbf{narrow the bounds} handed
    to each subtree so the ordering invariant cannot be broken:

\begin{minted}{python}
@st.composite
def bsts(draw, lo=-1000, hi=1000, max_depth=5):
    if max_depth <= 0 or lo > hi or draw(st.integers(0, 2)) == 0:
        return None
    v = draw(st.integers(min_value=lo, max_value=hi))
    return (draw(bsts(lo, v - 1, max_depth - 1)),
            v,
            draw(bsts(v + 1, hi, max_depth - 1)))
\end{minted}

    Everything in the left subtree is drawn from $[lo, v-1]$ and
    everything in the right from $[v+1, hi]$, so every comparison in
    the \Code{is_bst} predicate would succeed.

    The \Code{lo > hi} test matters: once the interval is empty there is
    no value left to draw, so the subtree must be \Code{None}.

    Another acceptable answer is to generate a sorted list of distinct
    values and build a tree from it. Note that if you're not careful
    in how you build the tree, this could only ever produce
    \emph{balanced} binary search trees, which is its own kind of
    bias.
  \end{solution}
\end{parts}
\newpage
% ========================================================
% Fuzzing
% ========================================================

\question[5] Explain how the coverage-guided fuzzers discussed in class
  work.
  \answerspace{20} %
  \begin{solution} %
    A grey-box fuzzer tests a program by repeatedly mutating inputs
    drawn from a corpus (``population'') of test inputs, seeded with an
    initial set. It adds lightweight instrumentation to the program so
    that each execution reports which branches were covered. If a
    generated input covers something new, it is added to the corpus and
    becomes a starting point for further mutation; otherwise it is
    discarded. The coverage feedback is what lets the fuzzer make
    progress into deeper code rather than re-testing the same shallow
    paths.
  \end{solution}

\question Suppose that a grey-box fuzzer is testing the following
  function, using the input \Code{"ab"} as its sole initial seed:

\begin{minted}{python}
def parse(s: str) -> int:
    if len(s) < 2:
        return 0
    if s[0] == '<':
        if s[-1] == '>':
            if 'x' in s:
                raise ValueError("boom")
            return 3
        return 2
    return 1
\end{minted}

  The fuzzer is equipped with the following mutation operations:
  \begin{itemize}
  \item delete a character at random,
  \item replace a character with another randomly chosen character, or
  \item insert a character at random.
  \end{itemize}

  At each iteration the fuzzer may apply either \textbf{one or two} of
  these mutations. For each of the following inputs, indicate whether
  the fuzzer could \textbf{produce} it from its initial seed; if so,
  indicate whether it would then \textbf{include} it in the population.
  (Indicate your answer by circling your choice.)

\begin{parts}

\part[2] \Code{"<b"}
  \begin{itemize}
  \item[ ] Produced from \Code{"ab"}: Yes \quad/\quad No %
    \hfill Added to the population: Yes \quad/\quad No
  \end{itemize}

\part[2] \Code{"xb"}
  \begin{itemize}
  \item[ ] Produced from \Code{"ab"}: Yes \quad/\quad No %
    \hfill Added to the population: Yes \quad/\quad No
  \end{itemize}

\part[2] \Code{"a"}
  \begin{itemize}
  \item[ ] Produced from \Code{"ab"}: Yes \quad/\quad No %
    \hfill Added to the population: Yes \quad/\quad No
  \end{itemize}

\part[2] \Code{"<>"}
  \begin{itemize}
  \item[ ] Produced from \Code{"ab"}: Yes \quad/\quad No %
    \hfill Added to the population: Yes \quad/\quad No
  \end{itemize}

\part[2] \Code{"<x>"}
  \begin{itemize}
  \item[ ] Produced from \Code{"ab"}: Yes \quad/\quad No %
    \hfill Added to the population: Yes \quad/\quad No
  \end{itemize}

\end{parts}

  % \begin{solutionorbox}[3in]
  %   The seed \Code{"ab"} takes the \Code{return 1} path, covering the
  %   two guards and that return. An input is added when it covers
  %   something the population has not covered before.

  %   \begin{center}\small
  %     \begin{tabular}{llll}
  %       & produced? & added? & \\ \hline
  %       \Code{"<b"}  & yes, 1 edit  & \textbf{yes} & \Code{a}$\to$\Code{<}; now enters the second \Code{if} and returns 2 \\
  %       \Code{"xb"}  & yes, 1 edit  & no           & \Code{a}$\to$\Code{x}; still returns 1, nothing new \\
  %       \Code{"a"}   & yes, 1 edit  & \textbf{yes} & delete \Code{b}; now trips the length guard and returns 0 \\
  %       \Code{"<>"}  & yes, 2 edits & \textbf{yes} & \Code{a}$\to$\Code{<}, \Code{b}$\to$\Code{>}; returns 3 \\
  %       \Code{"<x>"} & \textbf{no}  & ---          & three characters from a two-character seed needs more than two edits \\
  %     \end{tabular}
  %   \end{center}

  %   Note \Code{"a"}: it is \emph{shorter} than the seed, and it covers
  %   something new anyway. ``New coverage'' means reaching a line the
  %   population has not reached, not getting deeper into the program.
  % \end{solutionorbox}

\question
% --- from the grey-box quiz, part (c) --------------------
This function crashes when its input contains the substring
  \Code{"fuzz"}:

\begin{minted}{python}
def scan(s: str) -> None:
    matched = 0
    for c in s:
        if c == "fuzz"[matched]:
            matched += 1
            if matched == 4:
                crash()
        else:
            matched = 0
\end{minted}

  \begin{parts}
  \part[2] Suppose we have a grey-box fuzzer that keeps any input that
    covers a line none of its kept inputs covered. Would it keep
    \Code{"fu"} after already keeping \Code{"f"}? Explain in one or
    two sentences. %
  \begin{solutionorbox}[1.5in]
    No. The loop body is the same code on every iteration, so
    \Code{"fu"} executes exactly the same lines as \Code{"f"}, it just
    does so more times. Line coverage cannot tell that the input got
    further into the match.
  \end{solutionorbox}

% --- follow-on, from the same quiz -----------------------
\part[2] Given your answer above, why does AFL record a \emph{count}
  for each branch rather than a single bit saying whether the branch
  was taken? %
  \begin{solutionorbox}[1.5in]
    Because the count distinguishes looping once from looping twice,
    allowing the fuzzer to ``see'' that \Code{"fu"} exercises a
    different execution trace than \Code{"f"}.
  \end{solutionorbox}
\end{parts}

\newpage
% ========================================================
% Grammar-Based Fuzzing
% ========================================================

\question Consider the following grammar for arithmetic expressions:

\begin{center}
  \begin{tabular}{lcl}
    \texttt{<expr>}   & ::= & \texttt{<term> + <expr>}
                      $\mid$ \texttt{<term> - <expr>}
                      $\mid$ \texttt{<term>} \\
    \texttt{<term>}   & ::= & \texttt{<factor> * <term>}
                      $\mid$ \texttt{<factor> / <term>}
                      $\mid$ \texttt{<factor>} \\
    \texttt{<factor>} & ::= & \texttt{(<expr>)} $\mid$ \texttt{<digit>} \\
    \texttt{<digit>}  & ::= & \texttt{0} $\mid$ \texttt{1} $\mid$ \ldots
                      $\mid$ \texttt{9} \\
  \end{tabular}
\end{center}

\begin{parts}

\part[4] Give a derivation of the string \Code{2 * (3 + 4)} starting
  from \texttt{<expr>}. At each step, state which production you used.
  You may abbreviate runs of steps that only expand a
  \texttt{<factor>} into a single digit. %
  \answerspace{20} %
  \begin{solution} %
    One derivation (expansions of \texttt{<factor>} to a digit
    collapsed):

    \begin{center}\small
      \begin{tabular}{ll}
        \texttt{<expr>} & \\
        \texttt{<term>} & \texttt{<expr> ::= <term>} \\
        \texttt{<factor> * <term>} & \texttt{<term> ::= <factor> * <term>} \\
        \texttt{2 * <term>} & \texttt{<factor> ::= <digit> ::= 2} \\
        \texttt{2 * <factor>} & \texttt{<term> ::= <factor>} \\
        \texttt{2 * (<expr>)} & \texttt{<factor> ::= (<expr>)} \\
        \texttt{2 * (<term> + <expr>)} & \texttt{<expr> ::= <term> + <expr>} \\
        \texttt{2 * (3 + <expr>)} & \texttt{<term> ::= <factor> ::= <digit> ::= 3} \\
        \texttt{2 * (3 + <term>)} & \texttt{<expr> ::= <term>} \\
        \texttt{2 * (3 + 4)} & \texttt{<term> ::= <factor> ::= <digit> ::= 4} \\
      \end{tabular}
    \end{center}

  \end{solution}

\part[3] Suppose we instead have a grammar for URLs, and we use it to
  fuzz a program that fetches web pages. Even though every string the
  grammar produces is a syntactically valid URL, it turns out the
  program under test rejects almost half of them!  Can you think of a
  plausible explanation for why this might have happened? %
  \answerspace{10} %
  \begin{solution} %
    The grammar and the program disagree about what counts as a valid
    input. A typical URL grammar emits several schemes, e.g.,
    \Code{http}, \Code{https}, \Code{ftp}, \Code{ftps}, while a
    program that fetches web pages may only handle HTTP.

    The general takeaway is that a grammar is a \textbf{model} of the
    input format, and the model can be wrong: grammar-based generation
    buys validity \emph{with respect to the grammar we use}, which is
    only as good as our understanding of the format.
  \end{solution}

\end{parts}

\newpage
% ========================================================
% Minimization
% ========================================================

\question[5] Suppose we are fuzz testing a program and discover that the
  input string \Code{"axbxcxd"} causes the system under test to crash.
  Suppose that the delta debugging algorithm outputs the string
  \Code{"xx"}.

  Can you hypothesize what sorts of input strings trigger the failure?
  Using that hypothesis, explain how \Code{ddmin} went from the initial
  input to its final output by giving four intermediate inputs it might
  test, two of which induce the failure and two of which do not. List
  them in the order \Code{ddmin} would generate them.
  \answerspace{20} %
  \begin{solution} %
    A reasonable hypothesis is that the failure is triggered by any
    input containing at least two \Code{"x"}s, which is consistent with
    the minimal failing input \Code{"xx"}.

    Running \Code{ddmin} under that hypothesis, the first candidates it
    tests are:

    \begin{center}
      \begin{tabular}{ll}
        \Code{"axbx"} & \textcolor{red}{fails} \\
        \Code{"ax"}   & \textcolor{green}{passes} \\
        \Code{"bx"}   & \textcolor{green}{passes} \\
        \Code{"xbx"}  & \textcolor{red}{fails} \\
        \Code{"xx"}   & \textcolor{red}{fails} \\
      \end{tabular}
    \end{center}

    Any four inputs consistent with the hypothesis and with
    \Code{ddmin}'s subset/complement structure are acceptable.
  \end{solution}

% ========================================================
% Minimality
% ========================================================

\question Recall that \Code{ddmin} guarantees its output is
  \textbf{1-minimal}: removing any single remaining element makes the
  failure go away.

  Suppose a program reads a configuration file and crashes if
  \textbf{both} \Code{debug} and \Code{verbose} are enabled, \emph{or}
  if \Code{legacy} is enabled on its own. We run \Code{ddmin} on a
  configuration containing all three lines:

  \begin{center}
    \Code{["debug", "verbose", "legacy"]}
  \end{center}

\begin{parts}

\part[4] What does \Code{ddmin} return on this input? %
  \answerspace{22} %
  \begin{solution} %
    It returns \Code{["debug", "verbose"]}.

    \begin{center}\small
      \begin{tabular}{lll}
        & candidate & outcome \\ \hline
        subset     & \Code{["debug", "verbose"]} & \textbf{crashes} --- keep it, discard \Code{"legacy"} \\
        subset     & \Code{["debug"]}            & passes \\
        subset     & \Code{["verbose"]}          & passes \\
        complement & \Code{["verbose"]}          & passes \\
        complement & \Code{["debug"]}            & passes \\
      \end{tabular}
    \end{center}

    The answer \emph{is} 1-minimal: dropping either line stops the
    crash, so nothing can be removed.
  \end{solution}

\part[4] Would any of the other minimization strategies we have seen in
  class return a smaller result than \Code{ddmin} here? If so, why
  don't we always use that strategy to shrink failing inputs? %
  \answerspace{24} %
  \begin{solution} %
    Yes. The element-by-element shrinker that deletes one element at a
    time returns \Code{["legacy"]}, which is half the size of
    \Code{ddmin}'s output, in only three tests. Of course, we got
    lucky here: deleting left to right happens to remove
    \Code{"debug"} first on this example, but \Code{ddmin}'s
    partitioning never tries to do this.

    \textbf{Why we do not always use it:} its cost is linear in the
    size of the input, while \Code{ddmin} is logarithmic when the
    failure is localized.
  \end{solution}

\part[3] Can you think of a scenario where an input does \emph{not}
  cause the system under test to fail, but where we would still want to
  minimize it? %
  \answerspace{20} %
  \begin{solution} %
    The coverage-guided fuzzer from the project is a prime
    example. When an input is added to the population because it
    covered something new, the fuzzer trims it before spending
    mutations on it, because a smaller input has fewer positions to
    mutate.

    Other reasonable answers: minimizing a test case before adding it
    to a regression suite, or minimizing a reproducer for a
    performance problem rather than a crash.
  \end{solution}

\end{parts}

\newpage
% ========================================================
% Symbolic and Concolic Execution
% ========================================================

\question Consider the following function, which applies an adjustment
  to a test score, but caps the point total at 100 points:

\begin{minted}{python}
def grade(score: int, bonus: int) -> int:
    total: int = score
    if bonus > 0:
        total = total + bonus
    if total > 100:
        total = 100
    if total < 0:
        total = 0
    return total
\end{minted}

\begin{parts}

\part[6] Draw the symbolic execution tree that an engine would produce
  for \Code{grade}. Mark any unreachable nodes with an X; there is no
  need to expand the execution tree beneath that node. Assume the
  symbolic inputs are \Code{score} $\mapsto s$ and \Code{bonus}
  $\mapsto b$. %
  \answerspace{30} %
  \begin{solution} %
    Three \emph{sequential} conditionals, so the tree branches three
    times on every path and has $2^3 = 8$ leaves. Reading each path as
    a conjunction of its edge conditions:

    \begin{center}\small
      \begin{tabular}{cl}
        leaf & path condition \\ \hline
        1 & $b > 0 \wedge s + b > 100 \wedge 100 < 0$ \\
        2 & $b > 0 \wedge s + b > 100 \wedge \lnot(100 < 0)$ \\
        3 & $b > 0 \wedge \lnot(s + b > 100) \wedge s + b < 0$ \\
        4 & $b > 0 \wedge \lnot(s + b > 100) \wedge \lnot(s + b < 0)$ \\
        5 & $\lnot(b > 0) \wedge s > 100 \wedge 100 < 0$ \\
        6 & $\lnot(b > 0) \wedge s > 100 \wedge \lnot(100 < 0)$ \\
        7 & $\lnot(b > 0) \wedge \lnot(s > 100) \wedge s < 0$ \\
        8 & $\lnot(b > 0) \wedge \lnot(s > 100) \wedge \lnot(s < 0)$ \\
      \end{tabular}
    \end{center}

    Note the third condition: on any path that took the \Code{total >
      100} branch, \Code{total} has been \emph{reassigned} to
    \Code{100}, so the test \Code{total < 0} is really $100 <
    0$. Tracking that reassignment is the whole point of carrying a
    symbolic state.
  \end{solution}

\newpage

\part[4] How many test paths does this tree represent? %
  \answerspace{22} %
  \begin{solution} %

    Witnesses for each test path:

    \begin{center}\small
      \begin{tabular}{cll}
        leaf & input & returns \\ \hline
        2 & \Code{score=100, bonus=1} & 100 \\
        3 & \Code{score=-2, bonus=1}  & 0 \\
        4 & \Code{score=99, bonus=1}  & 100 \\
        6 & \Code{score=101, bonus=0} & 100 \\
        7 & \Code{score=-1, bonus=0}  & 0 \\
        8 & \Code{score=2, bonus=0}   & 2 \\
      \end{tabular}
    \end{center}
  \end{solution}

\part[3] Suppose we rewrite \Code{grade} so that the second
  conditional is \emph{nested} inside the first:

\begin{minted}{python}
    if bonus > 0:
        total = total + bonus
        if total > 100:
            total = 100
    if total < 0:
        total = 0
\end{minted}

  How many leaves does the tree have now? %
  \answerspace{20} %
  \begin{solution} %
    \textbf{Six} leaves, of which \textbf{five} are feasible.

    The \Code{total > 100} test is now only reached when \Code{bonus >
      0}; the \Code{bonus <= 0} side goes straight to the final
    conditional. That side now contributes 2 leaves rather than 4, and
    only one infeasible path survives (the one that clamps to 100 and
    then asks whether $100 < 0$).

    The lesson: $2^n$ is an \textbf{upper bound} on the number of paths
    through $n$ conditionals, reached only when the conditionals are
    independent. Nesting reduces the count --- and infeasibility
    reduces it further. Both work in our favour, but neither is enough
    to stop the number growing exponentially in a real program.
  \end{solution}

\end{parts}

\newpage

\question The following function sorts the first three elements of its
  argument in place:

\begin{minted}{python}
def bubble_sort3(arr: list) -> None:
    if len(arr) < 3:
        return
    for i in range(2):
        for j in range(2 - i):
            if arr[j] > arr[j + 1]:
                tmp: int = arr[j]
                arr[j] = arr[j + 1]
                arr[j + 1] = tmp
\end{minted}

\begin{parts}
\part[6] Provide three concrete test inputs for \Code{bubble_sort3}
  that induce distinct test paths. For each, give the final path
  constraint that concolic execution would produce. Write your answer
  as a sequence of conditions separated by $\land$. Assume the initial
  symbolic input has the form $arr \mapsto [a]$ or
  $arr \mapsto [a, b, c]$, i.e., use $a, b,$ and $c$ for the symbolic
  values stored at the first, second, and third indices. You may drop
  clauses from the path condition that are always true.
  \answerspace{26} %
  \begin{solution} %
    Many correct answers. One set:

    \begin{itemize}
    \item \Code{[1]}: \quad $\texttt{len(arr)} < 3$
    \item \Code{[1, 2, 3]}: \quad $\lnot(\texttt{len(arr)} < 3) \land
      \lnot(a > b) \land \lnot(b > c) \land \lnot(a > b)$
    \item  \Code{[3, 2, 1]}: \quad $\lnot(\texttt{len(arr)} < 3) \land
      (a > b) \land (a > c) \land (b > c)$
    \end{itemize}

    The key point: each comparison in the nested loops contributes one
    clause, in the order the loops execute them, and the swap changes
    which symbolic value sits at each index for later comparisons.

    On \Code{[3, 2, 1]} the array starts as $[a, b, c]$. The first
    comparison $a > b$ holds, so after the swap the array holds
    $[b, a, c]$. The second comparison is therefore \textbf{$a > c$},
    not $b > c$; it holds too, leaving $[b, c, a]$. The third
    comparison is then \textbf{$b > c$}.

  \end{solution}

\newpage

\part[4] For each of your three test inputs, give a next test input
  that a concolic execution engine might generate from it, or explain
  why it cannot generate anything. State the strategy you used.
  \answerspace{22} %
  \begin{solution} %
 The strategy is to negate the \textbf{last} clause of the path
    constraint and ask a solver for an input satisfying the result.

    \medskip
    From \Code{[1]}: negating $\texttt{len(arr)} < 3$ requires any
    array of length at least three --- e.g. \Code{[0, 0, 0]}.

    \medskip
    From \Code{[3, 2, 1]}: negating the final $b > c$ gives
    $(a > b) \land (a > c) \land \lnot(b > c)$ --- e.g.
    \Code{[-1, -2, -2]}.

    \medskip
    From \Code{[1, 2, 3]}: \textbf{there is no such input.} That path
    constraint ends with $\lnot(a > b)$, which is the same clause it
    began with, so negating it yields
    $\lnot(a > b) \land \lnot(b > c) \land (a > b)$ --- a
    contradiction. The solver returns \Code{UNSAT} and the engine must
    pick a different clause to negate; flipping the \emph{middle}
    clause instead gives \Code{[1, 3, 2]}.

    We'd give full credit for any answer that applies the
    negate-the-last-clause strategy correctly and recognises when it
    yields an unsatisfiable constraint.
  \end{solution}
\end{parts}

\question[1] Name one advantage that concolic execution has compared to
  coverage-guided grey-box fuzzing.
  \answerspace{12} %
  \begin{solution} %
    A concolic engine assumes complete access to the program under
    test, so it can \emph{guarantee} that each new test it generates
    explores a new path. A coverage-guided fuzzer offers no such
    guarantee --- it mutates at random and hopes.

    Any answer showing the student knows concolic execution is a
    white-box approach, and can contrast grey- and white-box testing,
    works here.
  \end{solution}

\question[1] Name one disadvantage that concolic execution has compared
  to coverage-guided grey-box fuzzing.
  \answerspace{12} %
  \begin{solution} %
    It requires access to the source code of the system under test, and
    it requires a solver to find new inputs --- which is slow, and
    which may fail outright when the path conditions are complicated
    (nonlinear arithmetic, hashes, and so on).
  \end{solution}

\end{questions}

\end{document}
